% from ega4-18-fr.tex
\oldpage[IV-4]{121}

$R^{-1}B=B_{\mathfrak n}$ est l'unique idéal maximal $\mathfrak nB_{\mathfrak n}$ de cet anneau, donc l'image réciproque de $\mathfrak p$ dans $B$ est égale à $\mathfrak n$. Mais d'autre part, si $\mathfrak q$ est l'image réciproque de $\mathfrak p$ dans $C$, on a $\mathfrak q\cap B=\mathfrak n$, et comme $\mathfrak n$ est maximal dans $B$ et $C$ une $B$-algèbre finie, $\mathfrak q$ est nécessairement un des idéaux maximaux de $C$; en outre, on a $u\notin\mathfrak q$ puisque $u\notin\mathfrak r$ et que $u\in B$, donc par hypothèse on a nécessairement $\mathfrak q=\mathfrak r$.

D'autre part, comme $B\subset C$, l'homomorphisme $g:R^{-1}B\to R^{-1}C$ est injectif ($0_{\mathrm I}$, 1.3.2); pour voir qu'il est surjectif, notons que $R^{-1}C$ est un $(R^{-1}B)$-module de type fini, et d'autre part que $\mathfrak mR^{-1}B$ est contenu dans l'idéal maximal de l'anneau local $R^{-1}B$; en vertu du lemme de Nakayama, il suffit de prouver que l'homomorphisme
\[
  R^{-1}B/\mathfrak mR^{-1}B\longrightarrow R^{-1}C/\mathfrak mR^{-1}C
\]
est surjectif. Mais en vertu de la première partie de la démonstration, $R^{-1}C/\mathfrak mR^{-1}C$ s'identifie à $C_{\mathfrak r}/\mathfrak mC_{\mathfrak r}$; par hypothèse cette $k$-algèbre est engendrée par l'image de $u$, et \emph{a fortiori} elle est égale à l'image de $R^{-1}B/\mathfrak mR^{-1}B$.

Considérons en second lieu le cas où $f$ est étale au point $x$. En remplaçant $X$ par un voisinage de $x$, on peut supposer que $X$ est un voisinage de $\mathfrak n$ dans $\Spec(B)$ (I, 7.2). Posons $B'=\Spec(A[T]/F.A[T])$ et soit $\mathfrak n'$ l'image réciproque de $\mathfrak n$ dans $B'$; comme l'image de $F'(T)$ dans $B'$ n'appartient pas à $\mathfrak n'$ par hypothèse, le morphisme $\Spec(B')\to\Spec(A)$ est étale au point $\mathfrak n'$ par (18.4.2, (ii)). Comme par hypothèse $\Spec(B)\to\Spec(A)$ est étale au point $\mathfrak n$, on en conclut (17.3.4) que $\Spec(B)\to\Spec(B')$ est étale au point $\mathfrak n$; mais comme ce morphisme est une immersion, il ne peut être étale en un point que si c'est un isomorphisme local en ce point (17.9.1), donc $B_{\mathfrak n}$ et $B'_{\mathfrak n'}$ sont isomorphes.

Enfin, supposons que $\sh{O}_{X,x}$ soit une $A$-algèbre formellement étale; avec les notations précédentes, il résulte de (17.1.5) que $B_{\mathfrak n}$ est une $B'_{\mathfrak n'}$-algèbre formellement étale; mais puisque l'homomorphisme $B'_{\mathfrak n'}\to B_{\mathfrak n}$ est surjectif, cela ne peut avoir lieu que si cet homomorphisme est \emph{bijectif} (0, 19.10.3, (i)). Ceci achève la démonstration de (18.4.6).

\begin{corollary}[18.4.7]
\label{IV.18.4.7}
Soient $f:X\to Y$ un morphisme localement de type fini, $x$ un point de $X$. Pour que $f$ soit formellement non ramifié au point $x$, il faut et il suffit qu'il existe un voisinage ouvert $U$ de $x$ tel que $f|U$ se factorise en $U\xrightarrow{j}X'\xrightarrow{h}Y$, où $h$ est un morphisme étale et $j$ une immersion fermée.
\end{corollary}

On peut évidemment se borner au cas où $Y=\Spec(R)$ est affine et $f$ de type fini. Si $A=\sh{O}_{Y,y}$, avec $y=f(x)$, la condition pour que $f$ soit formellement non ramifié au point $x$ équivaut, en vertu de (17.4.1.2), à dire que $\sh{O}_{X,x}$ est une $A$-algèbre formellement non ramifiée. S'il en est ainsi, on peut appliquer (18.4.6, (ii)); remplaçant au besoin $Y$ par un voisinage affine de $y$, on peut supposer (avec les notations de (18.4.6)) que le polynôme $F$ est l'image dans $A[T]$ d'un polynôme unitaire $G\in R[T]$. On pose alors $X'=\Spec(R[T]/G.R[T])$; soit $x'$ l'image du point $\mathfrak n$ de $\Spec(B)$ par le morphisme correspondant à l'homomorphisme composé $R[T]/G.R[T]\to A[T]/F.A[T]\to B$. Il résulte de (18.4.2) que le morphisme $h:X'\to Y$ correspondant à l'homomorphisme canonique $R\to R[T]/G.R[T]$ est étale au point $x'$, donc, en restreignant au besoin $X'$ et $Y$ à des voisinages ouverts de $x'$ et $y$ respectivement, on peut supposer $h$ étale. D'autre part, en vertu de (I, 6.5.1, (ii)) et de (I, 7.2), il résulte du fait que l'on a un homomorphisme local $\varphi:\sh{O}_{X',x'}\to\sh{O}_{X,x}$ que cet homomorphisme correspond à un morphisme

